\begin{lemma}[The energy as a semigroup integral]
For every $v \in L^2(\gamma_n)$,
\[
  \mathcal{E}(v)
    = \frac{1}{2\sqrt\pi}\int_0^\infty
        \bigl\langle (I - P_s)v,\, v\bigr\rangle\, \frac{ds}{s^{3/2}},
\]
an identity in $[0,+\infty]$: the integrand is nonnegative, and the right-hand
side is $+\infty$ exactly when $v$ lies outside the form domain $\mathcal{V}$,
in which case $\mathcal{E}(v) = +\infty$ by the convention of \noderef{Energy}.
Here $P_s$ is the semigroup of \noderef{OUSemigroup}.
\end{lemma}
\begin{proof}
Step 1 (spectral expansion of the integrand).  Fix $s > 0$.  By
\noderef{OUSemigroupSpectral} the Hermite coefficients of $P_s v$ are
$e^{-s|\alpha|}v_\alpha$, so by Parseval (\noderef{HermiteBasis})
\[
  \bigl\langle (I-P_s)v, v\bigr\rangle
    = \sum_{\alpha} \bigl(1 - e^{-s|\alpha|}\bigr)|v_\alpha|^2
    = \sum_{\alpha \neq 0} \bigl(1 - e^{-s|\alpha|}\bigr)|v_\alpha|^2 ,
\]
the index $\alpha = 0$ contributing $0$ because $|0| = 0$ by
\noderef{MultiIndexSum} and hence $1 - e^{0} = 0$.  Every summand is nonnegative,
since $0 \leq 1 - e^{-s|\alpha|}\leq 1$, and the sum lies between $0$ and
$\|v\|_2^2$.

Step 2 (Tonelli).  For each fixed $\alpha$ the map
$s \mapsto (1-e^{-s|\alpha|})s^{-3/2}|v_\alpha|^2$ is continuous and nonnegative
on $(0,\infty)$, hence measurable.  Tonelli's theorem, applied to counting
measure on $\mathbb{N}_0^n \setminus \{0\}$ and the measure
$s^{-3/2}\,ds$ on $(0,\infty)$, allows exchanging the sum and the integral
without any finiteness assumption, both sides being values in $[0,+\infty]$:
\[
  \frac{1}{2\sqrt\pi}\int_0^\infty
      \langle(I-P_s)v,v\rangle\,\frac{ds}{s^{3/2}}
    = \sum_{\alpha \neq 0}|v_\alpha|^2 \cdot
      \frac{1}{2\sqrt\pi}\int_0^\infty\bigl(1-e^{-s|\alpha|}\bigr)\frac{ds}{s^{3/2}} .
\]

Step 3 (evaluating each factor).  For $\alpha \neq 0$ we have
$|\alpha| \geq 1 > 0$ by \noderef{MultiIndexSum}, so
\noderef{ScalarFractional} applies with $\lambda = |\alpha| \geq 0$ and gives
\[
  \frac{1}{2\sqrt\pi}\int_0^\infty\bigl(1-e^{-s|\alpha|}\bigr)\frac{ds}{s^{3/2}}
    = |\alpha|^{1/2}.
\]
Substituting into Step 2,
\[
  \frac{1}{2\sqrt\pi}\int_0^\infty\langle(I-P_s)v,v\rangle\frac{ds}{s^{3/2}}
    = \sum_{\alpha\neq0}|\alpha|^{1/2}|v_\alpha|^2 ,
\]
and by \noderef{Energy} the right-hand side is exactly $\mathcal{E}(v)$, as an
identity in $[0,+\infty]$.  In particular the integral is finite precisely when
$\mathcal{E}(v) < \infty$, i.e. when $v \in \mathcal{V}$.
\end{proof}
